\chapter{DX-047 - Blocks App Binding Snippets}\label{record-143}

\noindent{\footnotesize\color{BookMuted}\path{docs/design/DX-047-blocks-app-binding-snippets.md}\par}\smallskip

\begin{RecordMeta}
id: DX-047
title: Blocks App Binding Snippets
status: active
lane: asap
priority: medium
github_issue: 342
legend: DX

\end{RecordMeta}

Legend: \href{https://github.com/flyingrobots/bijou/blob/1853ad902f6bf219326409e7f743a5a1d6517dbc/docs/legends/DX-developer-experience.md}{DX - Developer Experience}

\section{Decision Summary}\label{record-143-decision-summary}

DOGFOOD's Blocks docs should show how a block becomes application code, not only how a block is described as metadata. For \texttt{v\allowbreak{}7\allowbreak{}.\allowbreak{}2\allowbreak{}.\allowbreak{}0}, the bounded repair is to add an app-binding walkthrough centered on the existing \texttt{C\allowbreak{}o\allowbreak{}u\allowbreak{}n\allowbreak{}t\allowbreak{}e\allowbreak{}r\allowbreak{}D\allowbreak{}e\allowbreak{}m\allowbreak{}o\allowbreak{}B\allowbreak{}l\allowbreak{}o\allowbreak{}c\allowbreak{}k} fixture:

\begin{itemize}
\tightlist
\item
  show the block contract: model, config, data requirement, and command intents
\item
  show the application state/update path that owns counter state
\item
  show how keyboard actions become command-intent emissions
\item
  show how the app renders the block from model-owned state
\item
  prove the snippet is not stale by compiling a canonical source file and testing that the rendered docs fence matches it exactly
\end{itemize}

\section{Sponsored Human}\label{record-143-sponsored-human}

A release-video viewer needs to understand why Blocks are reusable product assemblies rather than screenshots or one-off demos. The Blocks page should let the presenter say ``this is how you wire one into your app'' without leaving DOGFOOD.

\section{Sponsored Agent}\label{record-143-sponsored-agent}

An agent needs a deterministic docs witness that can inspect the Blocks section and verify that the app-binding example names the real API pieces: block config, state ownership, command intents, update routing, render invocation, and lower-mode expectations.

\section{Hill}\label{record-143-hill}

After this cycle, a developer can open DOGFOOD's Blocks section and see a complete, real app-binding snippet for \texttt{C\allowbreak{}o\allowbreak{}u\allowbreak{}n\allowbreak{}t\allowbreak{}e\allowbreak{}r\allowbreak{}D\allowbreak{}e\allowbreak{}m\allowbreak{}o\allowbreak{}B\allowbreak{}l\allowbreak{}o\allowbreak{}c\allowbreak{}k}, including the state and command routing that make the block interactive.

\section{Problem}\label{record-143-problem}

Issue \#342 came from the \texttt{v\allowbreak{}7\allowbreak{}.\allowbreak{}1\allowbreak{}.\allowbreak{}0} release-video rehearsal. The Blocks section now has useful catalog and preview surfaces, but it still reads like a product taxonomy. It does not yet answer the practical developer question:

\begin{CodeBlock}
How do I bind this block to my application's state and commands?
\end{CodeBlock}

The current ``How to Make Your Own Blocks'' page shows block authoring metadata and a minimal \texttt{d\allowbreak{}e\allowbreak{}f\allowbreak{}i\allowbreak{}n\allowbreak{}e\allowbreak{}B\allowbreak{}l\allowbreak{}o\allowbreak{}c\allowbreak{}k\allowbreak{}(\allowbreak{})} shape. It does not show the application side of the contract: TEA state, update routing, render-time config, and command-intent handling.

\section{Scope}\label{record-143-scope}

\begin{itemize}
\tightlist
\item
  Add a focused app-binding section to the Blocks docs.
\item
  Use \texttt{C\allowbreak{}o\allowbreak{}u\allowbreak{}n\allowbreak{}t\allowbreak{}e\allowbreak{}r\allowbreak{}D\allowbreak{}e\allowbreak{}m\allowbreak{}o\allowbreak{}B\allowbreak{}l\allowbreak{}o\allowbreak{}c\allowbreak{}k} as the first example because it already proves bounded state, command intents, interaction keys, and lower-mode output in DOGFOOD.
\item
  Keep the code snippet small enough to read inside the terminal.
\item
  Back the snippet with deterministic tests so API names and command paths do not drift silently.
\item
  Update release-facing signposts and changelog entries for \#342.
\end{itemize}

\section{Non-Goals}\label{record-143-non-goals}

\begin{itemize}
\tightlist
\item
  Do not redesign BlockLab.
\item
  Do not build a code editor or tutorial runner.
\item
  Do not pull broad \#302 Runtime Graph / Scene IR product-contract work into this cycle.
\item
  Do not add a new production counter block. \texttt{C\allowbreak{}o\allowbreak{}u\allowbreak{}n\allowbreak{}t\allowbreak{}e\allowbreak{}r\allowbreak{}D\allowbreak{}e\allowbreak{}m\allowbreak{}o\allowbreak{}B\allowbreak{}l\allowbreak{}o\allowbreak{}c\allowbreak{}k} remains a DOGFOOD fixture.
\item
  Do not rewrite the standard block catalog.
\end{itemize}

\section{User Experience Contract}\label{record-143-user-experience-contract}

The Blocks docs should show both sides of a block contract:

\begin{itemize}
\tightlist
\item
  the reusable block owns metadata, data requirements, command intents, render contract, and lower-mode behavior
\item
  the application owns model state, update decisions, host data, and when to invoke the block
\end{itemize}

The snippet should be readable in DOGFOOD at normal release-video dimensions. It should not require a browser, external docs page, or source-code search to explain the basic integration path.

\section{Playback Questions}\label{record-143-playback-questions}

\begin{itemize}
\tightlist
\item
  Can a viewer see where counter state lives?
\item
  Can a viewer see which key emits which command intent?
\item
  Can a viewer see that the block render input is built from app-owned state?
\item
  Can a viewer see how lower modes stay meaningful when interaction is absent?
\item
  Can a test fail if the docs remove the binding path or let the snippet drift away from real APIs?
\end{itemize}

\section{Accessibility / Assistive Posture}\label{record-143-accessibility-assistive-posture}

The binding explanation must be useful as text. Code snippets and prose should not rely on color-only emphasis. The lower-mode note should explain what pipe and accessible output preserve from the interactive block.

\section{Localization / Directionality Posture}\label{record-143-localization-directionality-posture}

This cycle adds English-source documentation copy only. DOGFOOD already labels English-source Markdown under non-English locales; this cycle should preserve that behavior. No new right-to-left layout behavior is introduced.

\section{Agent Inspectability / Explainability Posture}\label{record-143-agent-inspectability-explainability-posture}

Agents should be able to inspect:

\begin{itemize}
\tightlist
\item
  the rendered Blocks docs text
\item
  the exact snippet source text
\item
  whether the snippet references real exported \texttt{C\allowbreak{}o\allowbreak{}u\allowbreak{}n\allowbreak{}t\allowbreak{}e\allowbreak{}r\allowbreak{}D\allowbreak{}e\allowbreak{}m\allowbreak{}o\allowbreak{}B\allowbreak{}l\allowbreak{}o\allowbreak{}c\allowbreak{}k} helpers
\item
  whether the snippet names command-intent routing and lower-mode output
\end{itemize}

The proof should not depend on screenshots.

\section{Linked Invariants}\label{record-143-linked-invariants}

\begin{itemize}
\tightlist
\item
  Blocks are product-level assemblies, not leaf visual components.
\item
  Application state owns business data; block render functions consume explicit inputs.
\item
  Command intents describe user intent; applications interpret and apply those intents.
\item
  DOGFOOD release claims must be demonstrable inside DOGFOOD.
\end{itemize}

\section{Implementation Outline}\label{record-143-implementation-outline}

\begin{enumerate}
\def\labelenumi{\arabic{enumi}.}
\tightlist
\item
  Add a reusable app-binding snippet source for \texttt{C\allowbreak{}o\allowbreak{}u\allowbreak{}n\allowbreak{}t\allowbreak{}e\allowbreak{}r\allowbreak{}D\allowbreak{}e\allowbreak{}m\allowbreak{}o\allowbreak{}B\allowbreak{}l\allowbreak{}o\allowbreak{}c\allowbreak{}k} at \texttt{e\allowbreak{}x\allowbreak{}a\allowbreak{}m\allowbreak{}p\allowbreak{}l\allowbreak{}e\allowbreak{}s\allowbreak{}/\allowbreak{}d\allowbreak{}o\allowbreak{}c\allowbreak{}s\allowbreak{}/\allowbreak{}c\allowbreak{}o\allowbreak{}n\allowbreak{}t\allowbreak{}e\allowbreak{}n\allowbreak{}t\allowbreak{}/\allowbreak{}b\allowbreak{}l\allowbreak{}o\allowbreak{}c\allowbreak{}k\allowbreak{}s\allowbreak{}-\allowbreak{}c\allowbreak{}o\allowbreak{}u\allowbreak{}n\allowbreak{}t\allowbreak{}e\allowbreak{}r\allowbreak{}-\allowbreak{}a\allowbreak{}p\allowbreak{}p\allowbreak{}-\allowbreak{}b\allowbreak{}i\allowbreak{}n\allowbreak{}d\allowbreak{}i\allowbreak{}n\allowbreak{}g\allowbreak{}.\allowbreak{}s\allowbreak{}n\allowbreak{}i\allowbreak{}p\allowbreak{}p\allowbreak{}e\allowbreak{}t\allowbreak{}.\allowbreak{}t\allowbreak{}s}.
\item
  Render that exact snippet in \texttt{e\allowbreak{}x\allowbreak{}a\allowbreak{}m\allowbreak{}p\allowbreak{}l\allowbreak{}e\allowbreak{}s\allowbreak{}/\allowbreak{}d\allowbreak{}o\allowbreak{}c\allowbreak{}s\allowbreak{}/\allowbreak{}c\allowbreak{}o\allowbreak{}n\allowbreak{}t\allowbreak{}e\allowbreak{}n\allowbreak{}t\allowbreak{}/\allowbreak{}b\allowbreak{}l\allowbreak{}o\allowbreak{}c\allowbreak{}k\allowbreak{}s\allowbreak{}-\allowbreak{}m\allowbreak{}a\allowbreak{}k\allowbreak{}e\allowbreak{}-\allowbreak{}y\allowbreak{}o\allowbreak{}u\allowbreak{}r\allowbreak{}-\allowbreak{}o\allowbreak{}w\allowbreak{}n\allowbreak{}.\allowbreak{}m\allowbreak{}d}; generated copies must derive from the canonical snippet source.
\item
  Add prose explaining app-owned state, command intent routing, render-time config, and lower-mode output.
\item
  Add focused tests that render the Blocks guide and assert the app-binding section is visible and legible.
\item
  Add compile and source-integrity tests so the snippet cannot drift from the canonical source file or the real helper API.
\item
  Update \texttt{d\allowbreak{}o\allowbreak{}c\allowbreak{}s\allowbreak{}/\allowbreak{}C\allowbreak{}H\allowbreak{}A\allowbreak{}N\allowbreak{}G\allowbreak{}E\allowbreak{}L\allowbreak{}O\allowbreak{}G\allowbreak{}.\allowbreak{}m\allowbreak{}d}, \texttt{d\allowbreak{}o\allowbreak{}c\allowbreak{}s\allowbreak{}/\allowbreak{}B\allowbreak{}E\allowbreak{}A\allowbreak{}R\allowbreak{}I\allowbreak{}N\allowbreak{}G\allowbreak{}.\allowbreak{}m\allowbreak{}d}, and \texttt{d\allowbreak{}o\allowbreak{}c\allowbreak{}s\allowbreak{}/\allowbreak{}R\allowbreak{}O\allowbreak{}A\allowbreak{}D\allowbreak{}M\allowbreak{}A\allowbreak{}P\allowbreak{}.\allowbreak{}m\allowbreak{}d}.
\end{enumerate}

\section{Tests To Write First}\label{record-143-tests-to-write-first}

\begin{itemize}
\tightlist
\item
  A DOGFOOD Blocks render test fails unless ``How to Make Your Own Blocks'' includes a visible \texttt{C\allowbreak{}o\allowbreak{}u\allowbreak{}n\allowbreak{}t\allowbreak{}e\allowbreak{}r\allowbreak{}D\allowbreak{}e\allowbreak{}m\allowbreak{}o\allowbreak{}B\allowbreak{}l\allowbreak{}o\allowbreak{}c\allowbreak{}k} app-binding section.
\item
  A snippet integrity test fails unless the documented code matches the canonical source file and references the real counter model, config, command-intent, update, and render helpers.
\item
  Existing CounterDemoBlock preview and intent-key tests continue to pass.
\end{itemize}

\section{Acceptance Criteria}\label{record-143-acceptance-criteria}

\begin{itemize}
\tightlist
\item
  Blocks docs include at least one complete app-binding snippet.
\item
  The snippet uses real Bijou / DOGFOOD APIs from a canonical source file and is covered by deterministic type and source-integrity proof.
\item
  Interactive block docs explain state ownership and command routing.
\item
  DOGFOOD renders the snippet legibly in the Blocks section.
\item
  The release/video script can explain ``how to use this in your app'' from the Blocks docs alone.
\item
  Issue \#342 is closed by the implementation PR.
\end{itemize}

\section{Retrospective}\label{record-143-retrospective}

To be completed when the PR lands.

